\chapter{Produits de colocation et leur mode de vente}\label{ch:nw:colocation-products-and-how-they-are-sold}

Les règles européennes exigent qu'une plateforme de négociation qui vend de la \omterm{def:nw:colocation-products-and-how-they-are-sold:colo}{colocation} donne à chaque utilisateur du même service accès à son réseau dans les
mêmes conditions, et la liste de ces conditions comprend, entre le refroidissement et l'accès aux données, la longueur du câble. Dans le New
Jersey, Nasdaq a indiqué à la SEC qu'il refait le câblage du campus de ses centres de données afin que chaque connexion d'opérateur télécom
atteigne ses clients sur une distance égale, et il a baptisé ce chantier «~Equalization Project~». Le client dont la baie
se trouve à côté du moteur d'appariement et le client situé à l'autre bout de la salle sont censés attendre les mêmes nanosecondes pour les mêmes
bits, et lorsque la salle ne rend pas les câbles égaux, une bobine de fibre s'en charge.

La partie~III de ce livre porte sur l'emplacement des bourses et sur ce qu'il en coûte d'en être proche. Ce chapitre décrit ce que vend le centre de données d'une plateforme
et à quelles conditions~: l'espace, l'alimentation électrique, le refroidissement et les \omterm{def:nw:colocation-products-and-how-they-are-sold:hands}{mains à distance}~; les \omterm{def:nw:colocation-products-and-how-they-are-sold:mmr}{interconnexions} et la \omterm{def:nw:colocation-products-and-how-they-are-sold:mmr}{salle d'interconnexion}~; les règles d'équité~; les
connexions au réseau propre à la plateforme~; et les grilles tarifaires, qui, aux États-Unis, sont des documents publics déposés.

\section{Ce que vend la colocation~: espace, alimentation électrique, refroidissement et mains à distance}

\begin{definition}[Colocation, hébergement de proximité]\label{def:nw:colocation-products-and-how-they-are-sold:colo}
La \emph{colocation}\index{colocation} désigne la location d'un espace, d'une alimentation électrique et d'un refroidissement pour les équipements d'un client, dans le centre de données qui abrite les
moteurs d'appariement d'une plateforme de négociation, avec des connexions au réseau de la plateforme. L'\emph{hébergement de proximité}\index{hébergement de proximité} (proximity hosting) désigne le même service proposé dans un
centre de données voisin que la plateforme n'exploite pas, dans le cadre d'un accord conclu avec la plateforme ou avec un tiers.
\end{definition}

\begin{definition}[Baie de colocation, cage de colocation, densité de puissance]\label{def:nw:colocation-products-and-how-they-are-sold:cabinet}
Une \emph{baie de colocation}\index{baie de colocation} (\omterm{def:nw:colocation-products-and-how-they-are-sold:colo}{colocation} cabinet) est une armoire fermable louée à un seul client, avec une enveloppe de puissance fixe. Une
\emph{cage de colocation}\index{cage de colocation} (\omterm{def:nw:colocation-products-and-how-they-are-sold:colo}{colocation} cage) est un espace grillagé de la salle, loué à un seul client, qui regroupe plusieurs baies et dont l'accès
est contrôlé à part. La \emph{densité de puissance}\index{densité de puissance} (power density) d'une baie est la puissance, en kilowatts, qu'elle est autorisée à consommer~; elle
fixe son prix tout autant que l'espace.
\end{definition}

Le centre de données d'une plateforme est un bâtiment industriel. Le centre de données Mahwah d'ICE, qui abrite les bourses NYSE, est un bâtiment de 398\,000
pieds carrés sur un site de 28~acres situé à 34~miles de Wall Street, avec 28~mégawatts de puissance électrique~; ses \omterm{def:nw:colocation-products-and-how-they-are-sold:cabinet}{baies de colocation} se déclinent en 4, 8 et
12~kilowatts, extensibles à 15, avec des baies partielles de 1 à 2~kilowatts, chacune alimentée par deux circuits électriques indépendants. Les baies se vendent par
\omterm{def:nw:colocation-products-and-how-they-are-sold:cabinet}{densité de puissance}, car la puissance et le refroidissement qui l'évacue sont les ressources rares du bâtiment~: le calcul du nombre de serveurs par baie du chapitre~8
est le côté client de la même contrainte.

\begin{definition}[Mains à distance]\label{def:nw:colocation-products-and-how-they-are-sold:hands}
Les \emph{mains à distance}\index{mains à distance} (remote hands) désigne le service de l'opérateur du centre de données qui réalise, sur demande du client, une intervention physique sur ses équipements
(remplacer une pièce, déplacer un câble, appuyer sur un bouton), facturée à la tâche ou à l'heure, afin que le client n'ait pas à envoyer
son propre personnel.
\end{definition}

\section{Interconnexions et salle d'interconnexion}

\begin{definition}[Interconnexion, salle d'interconnexion]\label{def:nw:colocation-products-and-how-they-are-sold:mmr}
Une \emph{interconnexion}\index{interconnexion} (cross-connect) est un câble, généralement une paire de fibres, installé et géré par l'opérateur du centre de données entre
deux points de son bâtiment~: la baie d'un client et le réseau de la plateforme, la baie d'un autre client, ou l'équipement d'un opérateur télécom. Une
\emph{salle d'interconnexion}\index{salle d'interconnexion} (meet-me room, MMR) est la zone du bâtiment où les opérateurs télécom terminent leurs réseaux et d'où
partent les interconnexions vers les clients.
\end{definition}

Toute connexion qui entre dans une salle de \omterm{def:nw:colocation-products-and-how-they-are-sold:colo}{colocation} ou la traverse est une \omterm{def:nw:colocation-products-and-how-they-are-sold:mmr}{interconnexion} dont quelqu'un paie les frais chaque mois. Les règles d'exploitation d'ICE pour
Mahwah sont typiques de ces règles~: les opérateurs ne peuvent s'installer que dans les \omterm{def:nw:colocation-products-and-how-they-are-sold:mmr}{salles d'interconnexion}, avec au plus 16~kilowatts chacun~; ils n'atteignent un client que
par une \omterm{def:nw:colocation-products-and-how-they-are-sold:mmr}{interconnexion}, pour laquelle des frais de connexion mensuels d'opérateur sont facturés~; et un client qui négocie sur la plateforme doit passer par le
réseau de la plateforme, et non par une \omterm{def:nw:colocation-products-and-how-they-are-sold:mmr}{interconnexion} directe vers ses systèmes. Les dépôts de Nasdaq décrivent la même structure à Carteret~: une cage d'opérateur, un
panneau de brassage, un câblage géré par la bourse jusqu'à un point de distribution, puis jusqu'à la baie du client, chaque élément «~codé par couleur, inventorié et
auditable~».

\begin{omfigure}
\begin{tikzpicture}[font=\footnotesize,b/.style={draw,rounded corners=2pt,minimum height=0.6cm,align=center,inner sep=3pt},>=Stealth]
\node[b,fill=omThm!15,minimum height=1.8cm] (me) at (0,0) {moteurs\\ d'appariement et\\ réseau de la\\ plateforme};
\node[b,fill=omDef!15] (c1) at (3.2,1.9) {baie proche};
\node[b,fill=omDef!15] (c2) at (6.2,1.9) {baie médiane};
\node[b,fill=omDef!15] (c3) at (9.2,1.9) {baie éloignée};
\draw[thick] (me.east |- 0,0.6) -- (3.2,0.6) -- (c1.south);
\draw[thick] (me.east |- 0,0.25) -- (6.2,0.25) -- (c2.south);
\draw[thick] (me.east |- 0,-0.1) -- (9.2,-0.1) -- (c3.south);
\foreach \r in {0.14,0.19,0.24} {\draw[omMeth,thick,fill=white] (2.2,0.6) circle (\r);}
\foreach \r in {0.14,0.19} {\draw[omMeth,thick,fill=white] (4.6,0.25) circle (\r);}
\node[font=\scriptsize,omMeth] at (2.2,1.05) {bobine longue};
\node[font=\scriptsize,omMeth] at (4.6,0.7) {bobine courte};
\node[b,fill=gray!12,minimum height=1.2cm] (mmr) at (0,-2.4) {salle d'interconnexion\\ (opérateurs)};
\node[b,fill=omDef!15] (cage) at (9.2,-1.4) {cage};
\draw[thick,gray] (mmr.north) -- node[left,font=\scriptsize]{liaisons des opérateurs} (me.south);
\draw[thick,gray] (mmr.east) -- (9.2,-2.4) node[pos=0.5,below,font=\scriptsize]{interconnexion depuis un opérateur} -- (cage.south);
\node[font=\scriptsize,align=center] at (6.2,-0.55) {chaque câble coupé à la longueur de la plus longue};
\end{tikzpicture}

\omcaption{Centre de données d'une plateforme de négociation, schématiquement. Les baies et cages des clients rejoignent le réseau de la plateforme par des câbles de longueur égale~;
plus une baie est proche, plus la fibre est enroulée sur son trajet. Les opérateurs terminent dans la \omterm{def:nw:colocation-products-and-how-they-are-sold:mmr}{salle d'interconnexion} et n'atteignent les clients que par
des \omterm{def:nw:colocation-products-and-how-they-are-sold:mmr}{interconnexions} gérées par l'opérateur.}\label{fig:nw:colocation-products-and-how-they-are-sold:hall}
\end{omfigure}

\section{Longueurs de câble égales et règles d'équité}

\begin{definition}[Égalisation de latence]\label{def:nw:colocation-products-and-how-they-are-sold:equal}
L'\emph{égalisation de latence}\index{egalisation de latence@égalisation de latence} (latency equalisation) désigne la pratique par laquelle une plateforme fait en sorte que le trajet de chaque client vers ses systèmes soit égal en
durée, quelle que soit la position du client dans le bâtiment, en coupant chaque connexion à la longueur de la plus longue (en enroulant le surplus)
ou en ajoutant un délai aux plus courtes.
\end{definition}

\begin{proposition}[Ce que coûte une bobine]\label{prop:nw:colocation-products-and-how-they-are-sold:coil}
Si le trajet le plus long vers le réseau de la plateforme mesure $L$ mètres de fibre d'indice de groupe $n_g$, une baie dont le trajet direct mesurerait $\ell$ mètres
attend $(L - \ell)\,n_g / c_0$ de plus que sans égalisation~: environ \qty{4.88}{\nano\second} par mètre pour une fibre standard.
\end{proposition}

\begin{proof}
Son câble est coupé à la longueur $L$~; la lumière parcourt les $L-\ell$ mètres supplémentaires à $c_0/n_g$.
\end{proof}

Une baie située à 5~mètres du moteur dans une salle dont le plus long parcours mesure 150~mètres abandonne donc \qty{707}{\nano\second}, le prix de
l'équité payé par le client qui aurait autrement gagné en se tenant tout près. La règle européenne fait de ce principe une obligation légale pour les plateformes
entrant dans son champ, en ajoutant que la plateforme doit surveiller «~toutes les connexions et mesures de latence~» et doit vendre chaque service de \omterm{def:nw:colocation-products-and-how-they-are-sold:colo}{colocation}
séparément plutôt qu'en bundles. Aux États-Unis, les services de \omterm{def:nw:colocation-products-and-how-they-are-sold:colo}{colocation} sont des services de bourse dont les tarifs sont déposés auprès de la SEC, et les
bourses décrivent leur égalisation dans ces dépôts. Ce que l'égalisation supprime, c'est la course vers la baie la plus proche~; ce qu'elle laisse subsister, ce sont toutes les autres
courses~: serveurs plus rapides, commutateurs plus rapides, meilleur code et, pour les sociétés qui négocient sur plusieurs plateformes, les routes entre les bâtiments des chapitres~10
à~14.

\section{Vitesses de port et produits de connectivité}

Une baie de client rejoint les systèmes de la plateforme par le réseau propre à celle-ci, au moyen de ports que la plateforme vend en plusieurs vitesses et
types~: ports de données de marché et ports de saisie d'ordres, d'un gigabit, de dix gigabits ou plus, un réseau standard et, sur certaines plateformes, un réseau à latence plus faible
à un prix plus élevé, ainsi que des connexions vers le site de reprise après sinistre de la plateforme (chapitre~28). Le prix d'une connexion est fixé par dépôt, comme
celui de la baie. Le produit le plus rapide de l'\cref{dat:nw:colocation-products-and-how-they-are-sold:fees}, la connexion d'une bourse d'options à 10~gigabits
à ultra-faible latence coûte dix fois sa connexion à 1~gigabit. Certaines plateformes vendent aussi une connectivité entre centres de données
(dont les liaisons sans fil du chapitre~14), et la règle des conditions égales s'applique à ce qu'elles vendent dans chaque produit.

\section{Barèmes de frais et achat d'un emplacement}

\begin{dated}{2026-09}{Frais de colocation et de connectivité, d'après les dépôts des bourses}\label{dat:nw:colocation-products-and-how-they-are-sold:fees}
Le dépôt de Nasdaq d'octobre 2024 fixe, pour sa salle d'extension NY11-4 à Carteret, des frais mensuels de 7\,230~USD pour une baie à très haute densité
(au-delà de 10 et jusqu'à \qty{15}{\kilo\watt}), des frais d'installation de baie de 5\,940~USD (baie comprise) et des frais d'installation d'alimentation de
3\,600~USD (monophasée) et 4\,560~USD (triphasée)~; il décrit ses frais mensuels de baie existants comme allant d'environ 475 à 916~USD
par kilowatt, et sa hausse de frais de novembre 2024 exclut les frais de NY11-4. Ses dépôts de 2026 étendent les mêmes frais d'installation à NY11-5
et listent des frais mensuels de circuit d'alimentation allant de 2\,640~USD (monophasé, 20~A, 240~V) à 12\,650.53~USD (triphasé, 32~A, 415~V). Le dépôt de MIAX Pearl
de mai 2026 a porté ses frais mensuels pour une connexion à 10~gigabits à ultra-faible latence vers ses installations principale et secondaire, de
13\,500 à 15\,000~USD, pour une connexion à 1~gigabit de 1\,400 à 1\,500~USD, et pour une connexion à 10~gigabits vers son site de reprise après sinistre
de 2\,750 à 3\,500~USD.
\end{dated}

\begin{method}[Acheter un emplacement en colocation]\label{met:nw:colocation-products-and-how-they-are-sold:buy}
\begin{enumerate}
\item Dimensionner d'abord l'alimentation (chapitre~8)~: les serveurs, commutateurs, dispositifs de synchronisation et de capture qui doivent se trouver dans le bâtiment, à leur consommation réelle.
\item Choisir les baies selon la \omterm{def:nw:colocation-products-and-how-they-are-sold:cabinet}{densité de puissance}~; une cage lorsque l'emplacement est grand ou que la société veut son propre contrôle d'accès.
\item Lister les connexions~: les ports de données de marché et de saisie d'ordres de la plateforme, les \omterm{def:nw:colocation-products-and-how-they-are-sold:mmr}{interconnexions} (cross-connects) vers les opérateurs pour les liaisons propres à la société, et les
  connexions de reprise après sinistre qu'exige la plateforme.
\item Chiffrer le tout d'après le barème actuel de la plateforme, frais ponctuels et mensuels, et amortir les frais ponctuels sur la durée du contrat.
\item Lire les règles d'exploitation de la plateforme~: ce qui peut être interconnecté avec quoi, qui peut entrer, les tarifs des \omterm{def:nw:colocation-products-and-how-they-are-sold:hands}{mains à distance} (remote hands), les délais de préavis.
\item Planifier les délais~: les baies, l'alimentation et les \omterm{def:nw:colocation-products-and-how-they-are-sold:mmr}{interconnexions} mettent plusieurs semaines à être livrées, et la certification d'une nouvelle connexion
  par la plateforme prend le sien.
\end{enumerate}
\end{method}

\begin{omfigure}
\begin{tikzpicture}
\begin{axis}[xbar stacked,width=11.5cm,height=5.8cm,bar width=9pt,xmin=0,xmax=440,
  ytick=data,yticklabels={une baie,deux baies,quatre baies,une connexion 1~Gb,deux 10~Gb ULL,deux 10~Gb ULL et secours},
  y dir=reverse,xlabel={coût annuel (milliers d'USD)},tick label style={font=\small},label style={font=\small},
  grid=major,grid style={gray!25},legend columns=2,legend style={font=\footnotesize,at={(0.5,-0.22)},anchor=north,draw=none,
  fill=none,/tikz/every even column/.append style={column sep=8pt}},table/col sep=comma]
\addplot[fill=omDef!55,draw=omDef] table[y expr=\coordindex,x=recurring_k]{figdata/networks/09-colocation-products-and-how-they-are-sold/bills.csv};
\addplot[fill=omThm!40,draw=omThm] table[y expr=\coordindex,x=one_time_k]{figdata/networks/09-colocation-products-and-how-they-are-sold/bills.csv};
\legend{frais mensuels,frais initiaux sur 36 mois}
\end{axis}
\end{tikzpicture}

\omcaption{Coût annuel de six emplacements d'après les frais publiés de l'\cref{dat:nw:colocation-products-and-how-they-are-sold:fees}~: baies Nasdaq NY11-4
(chacune avec installation et alimentation triphasée, la plus grande avec des unités de distribution d'alimentation) et connexions MIAX Pearl. Deux
connexions à ultra-faible latence coûtent à peu près autant que quatre baies. Données~: \texttt{fig\_colo.py} sur \texttt{firm.colobill}.}\label{fig:nw:colocation-products-and-how-they-are-sold:bills}
\end{omfigure}

Le graphique montre où va l'argent. Une baie chez Nasdaq coûte environ 90\,000~USD par an une fois son installation étalée sur trois ans~;
deux connexions à ultra-faible latence vers une bourse d'options coûtent 360\,000~USD par an, sans rien à installer. Le bâtiment est la
partie bon marché~; les connexions les plus rapides de la plateforme ne le sont pas, et une société les paie sur chaque plateforme où elle négocie.

\omcode{code/firm/colobill/firm_colobill.py}{56}{71}{Facture d'un emplacement~: frais mensuels, frais ponctuels amortis sur la durée du contrat, et la
vérification de date qui empêche un barème de devenir périmé.}\label{lst:nw:colocation-products-and-how-they-are-sold:bill}

\section{Tutoriel~: chiffrer un emplacement}\label{tut:nw:colocation-products-and-how-they-are-sold:tut}

\begin{tutorial}
\textbf{Objectif.} Transformer les frais publiés de deux plateformes en données, chiffrer six emplacements et celui d'une société, et calculer ce que coûte l'égalisation pour une
baie proche. \textbf{État final~:} \cref{fig:nw:colocation-products-and-how-they-are-sold:bills}, la facture de la société pour le problème, et le calcul
de la bobine.
\begin{tutsteps}
\item \textbf{Barèmes.} \texttt{firm.colobill.SCHEDULES} contient chaque frais avec la ligne du registre qui le source et la date de sa lecture.
\item \textbf{Emplacements.} \texttt{nw\_colo.FOOTPRINTS} liste les éléments et les quantités~; \texttt{bills()} les chiffre
  (\cref{lst:nw:colocation-products-and-how-they-are-sold:bill}) et \texttt{python fig\_colo.py} écrit les données du graphique.
\item \textbf{Ancienneté.} \texttt{stale(schedule, today)} liste les frais de plus de 60 jours~: la règle de l'encadré daté, appliquée dans le code.
\item \textbf{La bobine.} \texttt{coil(distances, longest)} évalue la \cref{prop:nw:colocation-products-and-how-they-are-sold:coil}.
\end{tutsteps}
\textbf{Pour aller plus loin.} Ajouter un troisième barème d'après le dépôt d'une autre plateforme et chiffrer le même emplacement là-bas~; porter la durée à 60
mois et voir combien pèsent alors peu les frais ponctuels.
\end{tutorial}

\section{Construction~: la facture de colocation}\label{bld:nw:colocation-products-and-how-they-are-sold:colobill}

\begin{build}
\bfield{Rôle} Les coûts de \omterm{def:nw:colocation-products-and-how-they-are-sold:colo}{colocation} et de connectivité de la société, calculés à partir des barèmes publiés, avec leur provenance et leurs dates, pour
le budget annuel du chapitre~29.
\bfield{Interface} \texttt{firm\_colobill}~: \texttt{Item(name, monthly, one\_time, unit, source, as\_of)},
\texttt{Schedule(venue, currency, items)}, \texttt{SCHEDULES}, \texttt{bill(schedule, footprint, term\_months)}, \texttt{stale(schedule, today, days)}, \texttt{break\_even},
\texttt{equalisation\_ns}.
\bfield{Règles} Chaque frais indique sa ligne du registre et sa date~; les frais ponctuels sont amortis sur une durée déclarée~; un élément inconnu est une erreur, et non un zéro~;
un frais de plus de 60 jours est signalé comme périmé.
\bfield{Tests d'acceptation} \texttt{code/firm/colobill/tests/}~: l'arithmétique d'une facture à la main, l'ancienneté à deux dates, le point d'équilibre, la bobine, et
un élément inconnu refusé.
\bfield{Pour aller plus loin} Lire un barème à partir du texte d'un dépôt~; ajouter des frais par paliers (un prix qui baisse avec le nombre de connexions).
\end{build}

\begin{omsources}
\item Commission Delegated Regulation (EU) 2017/573 (RTS~10), article~1.
\item SEC Releases 34-101267 (Nasdaq, SR-NASDAQ-2024-056) et 34-105010 (Nasdaq ISE, SR-ISE-2026-09)~; documents du Federal Register 2024-27759,
2026-06357 (Nasdaq) et 2026-10452 (MIAX Pearl).
\item ICE, \emph{ICE Global Network \& Colocation Technical Specifications} et \emph{Mahwah Operating Policies}.
\end{omsources}

\section{Exercices}

\begin{exercise}[$\star$]\label{exo:nw:colocation-products-and-how-they-are-sold:1}
Le plus long câble d'une salle mesure 120 mètres. Quel délai l'égalisation ajoute-t-elle à une baie dont le trajet direct serait de 30 mètres~?
\end{exercise}

\begin{exercise}[$\star$]\label{exo:nw:colocation-products-and-how-they-are-sold:2}
D'après l'\cref{dat:nw:colocation-products-and-how-they-are-sold:fees}, combien coûte par kilowatt et par mois une baie NY11-4 à très haute densité
à 10 puis à \qty{15}{\kilo\watt}~?
\end{exercise}

\begin{exercise}[$\star$]\label{exo:nw:colocation-products-and-how-they-are-sold:3}
Qu'a fait le dépôt de MIAX Pearl de mai 2026 au coût annuel de deux connexions à 10~gigabits à ultra-faible latence~?
\end{exercise}

\begin{exercise}[$\star\star$]\label{exo:nw:colocation-products-and-how-they-are-sold:4}
Pourquoi une plateforme vend-elle des baies selon la \omterm{def:nw:colocation-products-and-how-they-are-sold:cabinet}{densité de puissance} plutôt qu'en \omterm{def:nw:servers:ru}{unités de baie}~?
\end{exercise}

\begin{exercise}[$\star\star$]\label{exo:nw:colocation-products-and-how-they-are-sold:5}
Une société propose d'interconnecter directement sa baie à celle d'une autre société à Mahwah pour lui envoyer des ordres. Que disent les règles
d'exploitation de l'ICE, et pourquoi les plateformes écrivent-elles de telles règles~?
\end{exercise}

\begin{exercise}[$\star\star$]\label{exo:nw:colocation-products-and-how-they-are-sold:6}
L'égalisation supprime la course vers la baie la plus proche. Nommer trois autres courses qu'elle laisse subsister.
\end{exercise}

\begin{exercise}[$\star\star\star$]\label{exo:nw:colocation-products-and-how-they-are-sold:7}
\emph{Programmation.} Avec \texttt{firm.colobill}, chiffrer quatre baies NY11-4 avec installation, alimentation triphasée et PDU sur 36 et sur 60
mois. De combien la durée la plus longue réduit-elle le coût annuel~?
\end{exercise}

\begin{exercise}[$\star\star\star$]\label{exo:nw:colocation-products-and-how-they-are-sold:8}
\emph{Trouver la faille.} «~La plateforme égalise les longueurs de câble, donc chaque société colocalisée reçoit les données de marché au même instant.~»
\end{exercise}

\section{Problème~: La cage qui s'est payée elle-même}

\begin{problem}[{Problème du week-end --- un encombrement, sa facture et son seuil de rentabilité}]\label{pb:nw:colocation-products-and-how-they-are-sold:1}
Une société de market making colocalise deux Ultra High Density Cabinets dans la salle NY11-4 de Nasdaq (chacune avec installation et alimentation triphasée) et
achète deux connexions à 10~gigabits à ultra-faible latence vers une bourse d'options, aux tarifs de MIAX Pearl. Elle amortit les frais ponctuels sur 36 mois et
négocie 50\,000 contrats par jour, 20 jours par mois.

\medskip\noindent\textbf{Partie I --- La facture.}
\begin{enumerate}
\item Quels sont les frais mensuels des deux baies~?
\item Quels sont les frais ponctuels et leur amortissement mensuel~?
\item Combien coûtent les deux connexions par mois~?
\item Quel est le coût mensuel total de l'encombrement, et son coût annuel~?
\end{enumerate}

\medskip\noindent\textbf{Partie II --- Le seuil de rentabilité.}
\begin{enumerate}[resume]
\item Combien de contrats la société négocie-t-elle par mois~?
\item Que doit rapporter chaque contrat, net de tous les autres coûts, pour financer l'encombrement~?
\item Et si le volume était divisé par deux~?
\item Quelle ligne de la facture couperiez-vous en premier si c'était le cas~?
\end{enumerate}

\medskip\noindent\textbf{Partie III --- Équité et position.}
\begin{enumerate}[resume]
\item La baie de la société se trouve à 10 mètres du réseau de la plateforme, et le trajet le plus long fait 150 mètres. Que coûte l'égalisation à la société~?
\item Pourquoi la société ne peut-elle pas acheter un câble plus court~?
\item Qu'ajoute la règle européenne à cela~?
\item Que reçoit la société pour ses 15\,000~USD par mois et par connexion, que ne donne pas une connexion à 1~gigabit~?
\end{enumerate}

\medskip\noindent\textbf{Partie IV --- Le verdict.}
\begin{enumerate}[resume]
\item Énoncer le \emph{résultat nommé}~: le coût mensuel tout compris de l'encombrement et le seuil de rentabilité par contrat.
\item Sur trois ans, quelle part du total représentent les frais ponctuels~?
\item À quelle fréquence faut-il relire le barème des frais, et pourquoi~?
\item Que contient encore l'encombrement qu'les barèmes ne montrent pas~?
\item Comment la facture changerait-elle dans une plateforme où les connexions sont moins chères mais les baies plus coûteuses~?
\item Pourquoi est-ce la connexion, et non la baie, qui constitue la partie coûteuse~?
\item Que négocieriez-vous, et avec qui (le livre~16, chapitre~23)~?
\item En une phrase~: que achète un encombrement de \omterm{def:nw:colocation-products-and-how-they-are-sold:colo}{colocation}~?
\end{enumerate}
\end{problem}

\section{Questions d'entretien}

\begin{interviewq}[$\star$ \iqroles{développeur, trader}]\label{iq:nw:colocation-products-and-how-they-are-sold:1}
Qu'est-ce que la \omterm{def:nw:colocation-products-and-how-they-are-sold:colo}{colocation}, et que paie réellement une société~?
\end{interviewq}

\begin{interviewq}[$\star\star$ \iqroles{développeur}]\label{iq:nw:colocation-products-and-how-they-are-sold:2}
Qu'est-ce qu'un cross-connect et qu'est-ce qu'une meet-me room~? Comment le circuit d'un opérateur atteint-il votre baie~?
\end{interviewq}

\begin{interviewq}[$\star\star$ \iqroles{trader, développeur}]\label{iq:nw:colocation-products-and-how-they-are-sold:3}
Pourquoi les plateformes égalisent-elles les longueurs de câble, et quel avantage subsiste-t-il pour une société colocalisée~?
\end{interviewq}

\begin{interviewq}[$\star\star$ \iqroles{développeur}]\label{iq:nw:colocation-products-and-how-they-are-sold:4}
Vous emménagez dans le data center d'une plateforme dans trois mois. Que commandez-vous, et dans quel ordre~?
\end{interviewq}

\begin{interviewq}[$\star\star\star$ \iqroles{trader}]\label{iq:nw:colocation-products-and-how-they-are-sold:5}
Comment décideriez-vous si un encombrement de \omterm{def:nw:colocation-products-and-how-they-are-sold:colo}{colocation} sur une plateforme se rentabilise~?
\end{interviewq}

\begin{interviewq}[$\star\star$ \iqroles{développeur}]\label{iq:nw:colocation-products-and-how-they-are-sold:6}
Où trouvez-vous les frais de \omterm{def:nw:colocation-products-and-how-they-are-sold:colo}{colocation} d'une bourse américaine, et comment savez-vous qu'ils sont à jour~?
\end{interviewq}
